% =========================================================
% Chapter 3: Vectors
% POSN 1 Physics Preparation Notes
% Language: Thai
% Compiler: XeLaTeX
% =========================================================

\chapter{เวกเตอร์}

\begin{center}
    {\Large \textbf{Vectors}}\\[4pt]
    {\normalsize เอกสารเตรียมสอบคัดเลือก สอวน. ฟิสิกส์ ค่าย 1}\\[6pt]
    {\normalsize จัดทำโดย \textbf{Tames Thanakrit Damduan}}\\[2pt]
    {\footnotesize \textcopyright\ 2026 Tames Thanakrit Damduan. All rights reserved.}\\[8pt]
\end{center}

\section*{เป้าหมายของบทเรียน}

หลังจากเรียนบทนี้ นักเรียนควรทำสิ่งต่อไปนี้ได้

\begin{enumerate}
    \item เข้าใจความหมายของเวกเตอร์ ขนาด ทิศทาง และเวกเตอร์ตรงข้าม
    \item ใช้ระบบพิกัดฉากและระบบพิกัดเชิงขั้วในการระบุตำแหน่งได้
    \item แปลงพิกัดระหว่างรูป $(x,y)$ และรูป $(r,\theta)$ ได้
    \item บวก ลบ และคูณเวกเตอร์ด้วยจำนวนจริงได้
    \item แตกเวกเตอร์เป็นองค์ประกอบตามแกน $x$ และ $y$ ได้
    \item เขียนเวกเตอร์ในรูปเวกเตอร์หน่วย $\hat{i}$, $\hat{j}$, และ $\hat{k}$ ได้
    \item ใช้เวกเตอร์เป็นพื้นฐานในการเรียนแรง โมเมนตัม สนามไฟฟ้า และการเคลื่อนที่สองมิติได้
\end{enumerate}

\begin{tcolorbox}[title=แนวคิดสำคัญของบทนี้]
เวกเตอร์คือภาษาของปริมาณที่มีทั้งขนาดและทิศทาง จุดสำคัญที่สุดของบทนี้คือการแตกเวกเตอร์เป็นองค์ประกอบตามแกน เพราะหลังจากนั้นโจทย์สองมิติจะกลายเป็นโจทย์หนึ่งมิติแยกตามแกน
\end{tcolorbox}

\newpage{}

% =========================================================
\section{คณิตศาสตร์ก่อนเรียน: พิกัด มุม และตรีโกณมิติ}

\subsection{ทำไมต้องมีระบบพิกัด}

ในบทการเคลื่อนที่แนวตรง เราใช้แกนเดียว เช่น แกน $x$ เพื่อบอกตำแหน่ง แต่เมื่อจุดอยู่บนระนาบ เราต้องใช้ข้อมูลสองค่าเพื่อบอกตำแหน่งให้ครบ

ตัวอย่างเช่น จุดหนึ่งอยู่ทางขวา $3$ หน่วย และอยู่สูงขึ้น $4$ หน่วย จากจุดกำเนิด เราเขียนเป็น

\[
(x,y)=(3,4)
\]

ระบบนี้เรียกว่า ระบบพิกัดฉาก หรือ Cartesian coordinate system

\begin{tcolorbox}[title=นิยาม]
ระบบพิกัดฉาก คือ ระบบที่ใช้แกนตั้งฉากกัน เช่น แกน $x$ และแกน $y$ เพื่อบอกตำแหน่งของจุดในระนาบ
\end{tcolorbox}

\subsection{พิกัดฉาก}

ถ้าจุด $P$ มีพิกัด

\[
P=(x,y)
\]

หมายความว่า จากจุดกำเนิด ต้องเลื่อนไปตามแกน $x$ เป็นระยะ $x$ แล้วเลื่อนไปตามแกน $y$ เป็นระยะ $y$

ตัวอย่างเช่น

\[
P=(5,3)
\]

แปลว่า จุดนี้อยู่ทางขวา $5$ หน่วย และอยู่สูงขึ้น $3$ หน่วย

\[
Q=(-3,4)
\]

แปลว่า จุดนี้อยู่ทางซ้าย $3$ หน่วย และอยู่สูงขึ้น $4$ หน่วย

\subsection{พิกัดเชิงขั้ว}

บางครั้งการบอกตำแหน่งด้วยระยะจากจุดกำเนิดและมุมจะสะดวกกว่า เราเขียนในรูป

\[
(r,\theta)
\]

โดยที่

\[
r=\text{ระยะจากจุดกำเนิดถึงจุดนั้น}
\]

\[
\theta=\text{มุมที่วัดจากแกน }+x\text{ ทวนเข็มนาฬิกา}
\]

\begin{tcolorbox}[title=ข้อกำหนดในเอกสารนี้]
ถ้าไม่ได้บอกเป็นอย่างอื่น จะถือว่า $\theta$ วัดจากแกน $+x$ ทวนเข็มนาฬิกา
\end{tcolorbox}

\subsection{แปลงจากพิกัดเชิงขั้วเป็นพิกัดฉาก}

เริ่มจากจุดที่อยู่ห่างจากจุดกำเนิดเป็นระยะ $r$ และทำมุม $\theta$ กับแกน $+x$

จากนิยามของโคไซน์

\[
\cos\theta=\frac{\text{ด้านประชิดมุม}}{\text{ด้านตรงข้ามมุมฉาก}}
\]

ในรูปนี้ ด้านประชิดมุมคือ $x$ และด้านตรงข้ามมุมฉากคือ $r$

\[
\cos\theta=\frac{x}{r}
\]

จัดรูปได้

\[
x=r\cos\theta
\]

จากนิยามของไซน์

\[
\sin\theta=\frac{\text{ด้านตรงข้ามมุม}}{\text{ด้านตรงข้ามมุมฉาก}}
\]

ในรูปนี้ ด้านตรงข้ามมุมคือ $y$

\[
\sin\theta=\frac{y}{r}
\]

จัดรูปได้

\[
y=r\sin\theta
\]

\begin{tcolorbox}[title={สูตรแปลงพิกัดเชิงขั้วเป็นพิกัดฉาก}]
\[
x=r\cos\theta
\]

\[
y=r\sin\theta
\]
\end{tcolorbox}

\subsection{แปลงจากพิกัดฉากเป็นพิกัดเชิงขั้ว}

เริ่มจากจุดที่มีพิกัด

\[
(x,y)
\]

ระยะจากจุดกำเนิดถึงจุดนั้นหาได้จากทฤษฎีพีทาโกรัส

\[
r^2=x^2+y^2
\]

ดังนั้น

\[
r=\sqrt{x^2+y^2}
\]

มุมหาได้จากนิยามของแทนเจนต์

\[
\tan\theta=\frac{\text{ด้านตรงข้ามมุม}}{\text{ด้านประชิดมุม}}
\]

ในรูปนี้ ด้านตรงข้ามคือ $y$ และด้านประชิดคือ $x$

\[
\tan\theta=\frac{y}{x}
\]

\begin{tcolorbox}[title={สูตรแปลงพิกัดฉากเป็นพิกัดเชิงขั้ว}]
\[
r=\sqrt{x^2+y^2}
\]

\[
\tan\theta=\frac{y}{x}
\]
\end{tcolorbox}

\begin{tcolorbox}[title=จุดที่นักเรียนมักพลาด]
การใช้ $\theta=\tan^{-1}(y/x)$ ต้องระวังควอดแรนต์ เพราะเครื่องคิดเลขอาจให้มุมหลักที่ไม่ตรงกับทิศจริง ต้องดูเครื่องหมายของ $x$ และ $y$ ประกอบเสมอ
\end{tcolorbox}

\begin{examplebox}{ตัวอย่างที่ 1: แปลงพิกัดฉากเป็นพิกัดเชิงขั้ว}
จุด $P$ มีพิกัด

\[
P=(3,4)
\]

หาระยะจากจุดกำเนิด

\[
r=\sqrt{x^2+y^2}
\]

แทนค่า

\[
r=\sqrt{3^2+4^2}
\]

\[
r=5
\]

หามุม

\[
\tan\theta=\frac{y}{x}
\]

\[
\tan\theta=\frac{4}{3}
\]

ดังนั้น

\[
\theta=\tan^{-1}\left(\frac{4}{3}\right)
\]

จุดนี้อยู่ควอดแรนต์ที่ 1 เพราะ $x>0$ และ $y>0$
\end{examplebox}

\begin{exercisebox}{แบบฝึกหัดคณิตศาสตร์ก่อนเรียน}
\begin{enumerate}
    \item จงแปลงจุด $(6,8)$ ให้อยู่ในรูป $(r,\theta)$
    \item จงแปลงจุด $(r,\theta)=(10,30^\circ)$ ให้อยู่ในรูป $(x,y)$
    \item จุดหนึ่งมี $x=-3$ และ $y=4$ จงหา $r$ และระบุควอดแรนต์
    \item จุดหนึ่งมี $r=20$ และ $\theta=60^\circ$ จงหา $x$ และ $y$
\end{enumerate}
\end{exercisebox}

% =========================================================
\section{เวกเตอร์และการแทนเวกเตอร์}

\subsection{ความหมายของเวกเตอร์}

ปริมาณบางชนิดบอกแค่ขนาดอย่างเดียวไม่พอ ต้องบอกทิศทางด้วย เช่น การกระจัด ความเร็ว ความเร่ง และแรง ปริมาณแบบนี้เรียกว่า เวกเตอร์

\begin{tcolorbox}[title=นิยาม]
เวกเตอร์ คือ ปริมาณที่มีทั้งขนาดและทิศทาง
\end{tcolorbox}

เวกเตอร์มักเขียนด้วยลูกศรเหนือสัญลักษณ์ เช่น

\[
\vec{A}
\]

ขนาดของเวกเตอร์เขียนเป็น

\[
|\vec{A}|=A
\]

ขนาดของเวกเตอร์เป็นสเกลาร์ และต้องไม่ติดลบ

\[
A\geq 0
\]

\subsection{เวกเตอร์ที่เท่ากัน}

เวกเตอร์สองตัวจะเท่ากันเมื่อมี

\[
\text{ขนาดเท่ากัน}
\]

และ

\[
\text{ทิศทางเดียวกัน}
\]

ตำแหน่งที่วาดไม่สำคัญ เพราะเวกเตอร์สามารถเลื่อนขนานไปที่อื่นได้โดยยังเป็นเวกเตอร์เดิม ตราบใดที่ขนาดและทิศทางไม่เปลี่ยน

\begin{tcolorbox}[title=ข้อควรจำ]
เวกเตอร์ไม่จำเป็นต้องเริ่มที่จุดกำเนิด สิ่งที่กำหนดเวกเตอร์คือขนาดและทิศทาง ไม่ใช่ตำแหน่งที่วาด
\end{tcolorbox}

\subsection{เวกเตอร์ตรงข้ามและเวกเตอร์ศูนย์}

เวกเตอร์ตรงข้ามของ $\vec{A}$ เขียนเป็น

\[
-\vec{A}
\]

มีขนาดเท่ากับ $\vec{A}$ แต่มีทิศตรงข้าม

ดังนั้น

\[
\vec{A}+(-\vec{A})=\vec{0}
\]

โดย $\vec{0}$ เรียกว่า เวกเตอร์ศูนย์ มีขนาดเป็นศูนย์ และไม่มีทิศทางเฉพาะ

\begin{examplebox}{ตัวอย่างที่ 2: เวกเตอร์ตรงข้าม}
ถ้า $\vec{A}$ มีขนาด $12\,\si{N}$ ชี้ไปทางขวา แล้ว $-\vec{A}$ จะมีขนาดเท่ากับ

\[
12\,\si{N}
\]

แต่ชี้ไปทางซ้าย
\end{examplebox}

\begin{exercisebox}{แบบฝึกหัด 3.1}
\begin{enumerate}
    \item จงจำแนกปริมาณต่อไปนี้ว่าเป็นสเกลาร์หรือเวกเตอร์: แรง, พลังงาน, ความเร็ว, อัตราเร็ว, สนามไฟฟ้า, เวลา
    \item เวกเตอร์สองตัวมีขนาดเท่ากัน แต่ตัวหนึ่งชี้ทิศตะวันออก อีกตัวหนึ่งชี้ทิศตะวันตก เวกเตอร์สองตัวนี้เท่ากันหรือไม่
    \item ถ้า $\vec{A}$ มีขนาด $12\,\si{N}$ ไปทางขวา จงอธิบายความหมายของ $-\vec{A}$
    \item จงอธิบายว่าเหตุใดการกระจัดจึงเป็นเวกเตอร์ แต่ระยะทางเป็นสเกลาร์
\end{enumerate}
\end{exercisebox}

% =========================================================
\section{การบวก ลบ และคูณเวกเตอร์}

\subsection{ทำไมการบวกเวกเตอร์ไม่เหมือนการบวกจำนวนธรรมดา}

นักเรียนเดินจากจุดเริ่มต้นไปทางตะวันออก $3\,\si{m}$ แล้วเดินไปทางเหนือ $4\,\si{m}$

ระยะทางรวมคือ

\[
3+4=7\,\si{m}
\]

แต่การกระจัดลัพธ์ไม่ใช่ $7\,\si{m}$ เพราะสองทิศตั้งฉากกัน การกระจัดลัพธ์คือเส้นตรงจากจุดเริ่มต้นถึงจุดสุดท้าย

\[
R=\sqrt{3^2+4^2}=5\,\si{m}
\]

ดังนั้น การบวกเวกเตอร์ต้องคำนึงถึงทิศทาง ไม่ใช่บวกเฉพาะขนาด

\subsection{วิธีหัวต่อหาง}

ต้องการหา

\[
\vec{R}=\vec{A}+\vec{B}
\]

ให้คิดดังนี้

\begin{enumerate}
    \item วาดเวกเตอร์ $\vec{A}$
    \item นำหางของ $\vec{B}$ ไปต่อที่หัวของ $\vec{A}$
    \item เวกเตอร์ลัพธ์ $\vec{R}$ คือเวกเตอร์จากหางของ $\vec{A}$ ไปยังหัวของ $\vec{B}$
\end{enumerate}

\begin{tcolorbox}[title=กฎการบวกเวกเตอร์]
\[
\vec{A}+\vec{B}=\vec{B}+\vec{A}
\]

\[
\vec{A}+(\vec{B}+\vec{C})=(\vec{A}+\vec{B})+\vec{C}
\]
\end{tcolorbox}

\subsection{การลบเวกเตอร์}

การลบเวกเตอร์คือการบวกเวกเตอร์ตรงข้าม

\[
\vec{A}-\vec{B}=\vec{A}+(-\vec{B})
\]

ดังนั้นถ้าต้องการหา $\vec{A}-\vec{B}$ ให้กลับทิศของ $\vec{B}$ ก่อน แล้วนำไปบวกกับ $\vec{A}$

\subsection{การคูณเวกเตอร์ด้วยจำนวนจริง}

ถ้า $c$ เป็นจำนวนจริง การคูณ

\[
c\vec{A}
\]

มีผลดังนี้

\begin{enumerate}
    \item ขนาดกลายเป็น $|c|A$
    \item ถ้า $c>0$ ทิศทางเหมือนเดิม
    \item ถ้า $c<0$ ทิศทางตรงข้าม
    \item ถ้า $c=0$ ได้เวกเตอร์ศูนย์
\end{enumerate}

\begin{examplebox}{ตัวอย่างที่ 3: บวกเวกเตอร์ตั้งฉาก}
นักเรียนเดินไปทางตะวันออก $6\,\si{m}$ แล้วเดินไปทางเหนือ $8\,\si{m}$ จงหาขนาดของการกระจัดลัพธ์

เพราะสองทิศตั้งฉากกัน ใช้พีทาโกรัส

\[
R=\sqrt{6^2+8^2}
\]

\[
R=10\,\si{m}
\]

ถ้าต้องการหามุมกับทิศตะวันออก

\[
\tan\theta=\frac{8}{6}
\]

\[
\theta=\tan^{-1}\left(\frac{4}{3}\right)
\]
\end{examplebox}

\begin{exercisebox}{แบบฝึกหัด 3.2}
\begin{enumerate}
    \item นักเรียนเดินไปทางตะวันออก $5\,\si{m}$ แล้วเดินไปทางเหนือ $12\,\si{m}$ จงหาขนาดของการกระจัดลัพธ์
    \item เวกเตอร์ $\vec{A}$ มีขนาด $10\,\si{N}$ ไปทางขวา และ $\vec{B}$ มีขนาด $4\,\si{N}$ ไปทางซ้าย จงหา $\vec{A}+\vec{B}$
    \item ถ้า $\vec{A}$ ไปทางเหนือ และ $\vec{B}$ ไปทางตะวันออก จงอธิบายทิศของ $\vec{A}+\vec{B}$
    \item ถ้า $\vec{A}$ ชี้ขึ้นบน จงอธิบายทิศของ $-2\vec{A}$
\end{enumerate}
\end{exercisebox}

% =========================================================
\section{องค์ประกอบของเวกเตอร์}

\subsection{ทำไมต้องแตกเวกเตอร์}

การวาดรูปช่วยให้เห็นทิศทาง แต่โจทย์สอบต้องการคำตอบที่แม่นยำ วิธีที่เป็นระบบคือแตกเวกเตอร์เป็นองค์ประกอบตามแกน

แนวคิดคือ แทนเวกเตอร์เอียงหนึ่งตัวด้วยเวกเตอร์สองตัวที่ตั้งฉากกัน

\[
\vec{A}=\vec{A}_x+\vec{A}_y
\]

โดย $\vec{A}_x$ อยู่ตามแกน $x$ และ $\vec{A}_y$ อยู่ตามแกน $y$

\subsection{หาองค์ประกอบจากขนาดและมุม}

ให้เวกเตอร์ $\vec{A}$ มีขนาด $A$ และทำมุม $\theta$ กับแกน $+x$

จากนิยามของโคไซน์

\[
\cos\theta=\frac{A_x}{A}
\]

ดังนั้น

\[
A_x=A\cos\theta
\]

จากนิยามของไซน์

\[
\sin\theta=\frac{A_y}{A}
\]

ดังนั้น

\[
A_y=A\sin\theta
\]

\begin{tcolorbox}[title={สูตรแตกเวกเตอร์เมื่อมุมวัดจากแกนบวก x}]
\[
A_x=A\cos\theta
\]

\[
A_y=A\sin\theta
\]
\end{tcolorbox}

\begin{tcolorbox}[title=ระวัง]
สูตร $A_x=A\cos\theta$ และ $A_y=A\sin\theta$ ใช้ตรง ๆ เมื่อ $\theta$ วัดจากแกน $+x$ เท่านั้น ถ้ามุมวัดจากแกนอื่น ต้องดูจากรูปว่าส่วนใดเป็นด้านประชิดและส่วนใดเป็นด้านตรงข้าม
\end{tcolorbox}

\subsection{หาเวกเตอร์จากองค์ประกอบ}

ถ้ารู้ $A_x$ และ $A_y$ ขนาดของเวกเตอร์หาได้จากพีทาโกรัส

\[
A^2=A_x^2+A_y^2
\]

ดังนั้น

\[
A=\sqrt{A_x^2+A_y^2}
\]

มุมหาได้จาก

\[
\tan\theta=\frac{A_y}{A_x}
\]

แต่ต้องตรวจควอดแรนต์จากเครื่องหมายของ $A_x$ และ $A_y$

\begin{tcolorbox}[title=เครื่องหมายขององค์ประกอบ]
\[
\begin{array}{c|c|c}
\text{ควอดแรนต์} & A_x & A_y\\
\hline
\text{I} & + & +\\
\text{II} & - & +\\
\text{III} & - & -\\
\text{IV} & + & -
\end{array}
\]
\end{tcolorbox}

\begin{examplebox}{ตัวอย่างที่ 4: แตกเวกเตอร์เป็นองค์ประกอบ}
แรง $\vec{F}$ มีขนาด $50\,\si{N}$ ทำมุม $37^\circ$ กับแกน $+x$ จงหา $F_x$ และ $F_y$

เริ่มจากสูตรเมื่อมุมวัดจากแกน $+x$

\[
F_x=F\cos\theta
\]

\[
F_y=F\sin\theta
\]

แทนค่า

\[
F_x=50\cos37^\circ
\]

\[
F_y=50\sin37^\circ
\]

ใช้ค่าประมาณ

\[
\cos37^\circ\approx \frac{4}{5}
\]

\[
\sin37^\circ\approx \frac{3}{5}
\]

จึงได้

\[
F_x=50\left(\frac{4}{5}\right)=40\,\si{N}
\]

\[
F_y=50\left(\frac{3}{5}\right)=30\,\si{N}
\]
\end{examplebox}

\begin{examplebox}{ตัวอย่างที่ 5: มุมวัดจากแกน y}
แรง $\vec{F}$ มีขนาด $20\,\si{N}$ ทำมุม $30^\circ$ กับแกน $+y$ ไปทางขวา จงหา $F_x$ และ $F_y$

ครั้งนี้มุมไม่ได้วัดจากแกน $+x$ จึงต้องดูสามเหลี่ยม

องค์ประกอบตามแกน $y$ เป็นด้านประชิดมุม

\[
F_y=F\cos30^\circ
\]

องค์ประกอบตามแกน $x$ เป็นด้านตรงข้ามมุม

\[
F_x=F\sin30^\circ
\]

ดังนั้น

\[
F_x=20\sin30^\circ=10\,\si{N}
\]

\[
F_y=20\cos30^\circ=10\sqrt{3}\,\si{N}
\]
\end{examplebox}

\begin{exercisebox}{แบบฝึกหัด 3.3}
\begin{enumerate}
    \item เวกเตอร์ขนาด $10$ ทำมุม $30^\circ$ กับแกน $+x$ จงหา $A_x$ และ $A_y$
    \item เวกเตอร์หนึ่งมี $A_x=6$ และ $A_y=8$ จงหาขนาดของเวกเตอร์
    \item เวกเตอร์หนึ่งมี $A_x=-3$ และ $A_y=4$ จงหาขนาดและระบุควอดแรนต์
    \item แรง $100\,\si{N}$ ทำมุม $60^\circ$ กับแนวดิ่งไปทางขวา จงหาองค์ประกอบตามแนวนอนและแนวดิ่ง
    \item เวกเตอร์หนึ่งมีขนาด $20$ และชี้ลงทางขวาทำมุม $30^\circ$ ใต้แกน $+x$ จงหาองค์ประกอบ $x$ และ $y$
\end{enumerate}
\end{exercisebox}

% =========================================================
\section{เวกเตอร์หน่วยและการเขียนเวกเตอร์ด้วยองค์ประกอบ}

\subsection{เวกเตอร์หน่วย}

เวกเตอร์หน่วยคือเวกเตอร์ที่มีขนาดเท่ากับ $1$ และใช้บอกทิศทาง

\[
|\hat{i}|=|\hat{j}|=|\hat{k}|=1
\]

โดย

\[
\hat{i}=\text{เวกเตอร์หน่วยตามแกน }+x
\]

\[
\hat{j}=\text{เวกเตอร์หน่วยตามแกน }+y
\]

\[
\hat{k}=\text{เวกเตอร์หน่วยตามแกน }+z
\]

\begin{tcolorbox}[title=ความหมายของเวกเตอร์หน่วย]
เวกเตอร์หน่วยใช้บอกทิศทางขององค์ประกอบ ไม่ได้ใช้บอกขนาดของปริมาณทางฟิสิกส์
\end{tcolorbox}

\subsection{เขียนเวกเตอร์ในรูปเวกเตอร์หน่วย}

ถ้าเวกเตอร์ $\vec{A}$ มีองค์ประกอบ $A_x$ และ $A_y$ เราเขียนได้ว่า

\[
\vec{A}=A_x\hat{i}+A_y\hat{j}
\]

ในสามมิติ เขียนเป็น

\[
\vec{A}=A_x\hat{i}+A_y\hat{j}+A_z\hat{k}
\]

ขนาดของเวกเตอร์ในสามมิติคือ

\[
A=\sqrt{A_x^2+A_y^2+A_z^2}
\]

\begin{examplebox}{ตัวอย่างที่ 6: เวกเตอร์หน่วย}
เวกเตอร์หนึ่งมีองค์ประกอบ $x$ เท่ากับ $3$ และองค์ประกอบ $y$ เท่ากับ $-4$

เขียนได้ว่า

\[
\vec{A}=3\hat{i}-4\hat{j}
\]

ขนาดคือ

\[
A=\sqrt{3^2+(-4)^2}=5
\]

เวกเตอร์นี้อยู่ควอดแรนต์ที่ IV เพราะ $A_x>0$ และ $A_y<0$
\end{examplebox}

\subsection{เวกเตอร์ตำแหน่ง}

ถ้าจุด $P$ อยู่ที่ตำแหน่ง

\[
(x,y)
\]

เวกเตอร์จากจุดกำเนิดไปยังจุด $P$ เรียกว่า เวกเตอร์ตำแหน่ง

\[
\vec{r}=x\hat{i}+y\hat{j}
\]

ในสามมิติ

\[
\vec{r}=x\hat{i}+y\hat{j}+z\hat{k}
\]

\begin{tcolorbox}[title=เชื่อมกับบทถัดไป]
ในบทการเคลื่อนที่สองมิติ ตำแหน่งของอนุภาคจะเขียนเป็นเวกเตอร์ตำแหน่ง $\vec{r}(t)$ และความเร็วจะได้จากการเปลี่ยนแปลงของเวกเตอร์ตำแหน่งตามเวลา
\end{tcolorbox}

\begin{exercisebox}{แบบฝึกหัด 3.4}
\begin{enumerate}
    \item จงเขียนเวกเตอร์ที่มี $A_x=5$ และ $A_y=12$ ในรูปเวกเตอร์หน่วย
    \item จงหาขนาดของ $\vec{A}=8\hat{i}-6\hat{j}$
    \item จุด $P$ อยู่ที่ $(4,-7)$ จงเขียนเวกเตอร์ตำแหน่ง $\vec{r}$
    \item เวกเตอร์ $\vec{B}=-3\hat{i}-4\hat{j}$ อยู่ควอดแรนต์ใด และมีขนาดเท่าไร
    \item เวกเตอร์หนึ่งในสามมิติคือ $\vec{C}=2\hat{i}-3\hat{j}+6\hat{k}$ จงหาขนาดของเวกเตอร์นี้
\end{enumerate}
\end{exercisebox}

% =========================================================
\section{การบวกและลบเวกเตอร์ด้วยองค์ประกอบ}

\subsection{แนวคิดหลัก}

ถ้าต้องการบวกเวกเตอร์ให้แม่นยำ ให้แตกทุกเวกเตอร์เป็นองค์ประกอบก่อน แล้วบวกทีละแกน

ให้

\[
\vec{A}=A_x\hat{i}+A_y\hat{j}
\]

และ

\[
\vec{B}=B_x\hat{i}+B_y\hat{j}
\]

ต้องการหา

\[
\vec{R}=\vec{A}+\vec{B}
\]

แทนลงไป

\[
\vec{R}=(A_x\hat{i}+A_y\hat{j})+(B_x\hat{i}+B_y\hat{j})
\]

จัดกลุ่มแกนเดียวกัน

\[
\vec{R}=(A_x+B_x)\hat{i}+(A_y+B_y)\hat{j}
\]

ดังนั้น

\[
R_x=A_x+B_x
\]

\[
R_y=A_y+B_y
\]

\begin{tcolorbox}[title=สูตรบวกเวกเตอร์ด้วยองค์ประกอบ]
\[
R_x=A_x+B_x
\]

\[
R_y=A_y+B_y
\]

\[
\vec{R}=R_x\hat{i}+R_y\hat{j}
\]
\end{tcolorbox}

\subsection{หาขนาดและทิศของเวกเตอร์ลัพธ์}

หลังจากได้ $R_x$ และ $R_y$ แล้ว

\[
R=\sqrt{R_x^2+R_y^2}
\]

และ

\[
\tan\theta=\frac{R_y}{R_x}
\]

ต้องตรวจควอดแรนต์จากเครื่องหมายของ $R_x$ และ $R_y$ เสมอ

\subsection{การลบเวกเตอร์ด้วยองค์ประกอบ}

ให้

\[
\vec{D}=\vec{A}-\vec{B}
\]

จากนิยามการลบเวกเตอร์

\[
\vec{A}-\vec{B}=\vec{A}+(-\vec{B})
\]

ดังนั้นในรูปองค์ประกอบ

\[
D_x=A_x-B_x
\]

\[
D_y=A_y-B_y
\]

\[
\vec{D}=(A_x-B_x)\hat{i}+(A_y-B_y)\hat{j}
\]

\begin{examplebox}{ตัวอย่างที่ 7: บวกเวกเตอร์ด้วยองค์ประกอบ}
กำหนดให้

\[
\vec{A}=3\hat{i}+4\hat{j}
\]

และ

\[
\vec{B}=5\hat{i}-2\hat{j}
\]

หา

\[
\vec{R}=\vec{A}+\vec{B}
\]

บวกองค์ประกอบตามแกน $x$

\[
R_x=3+5=8
\]

บวกองค์ประกอบตามแกน $y$

\[
R_y=4+(-2)=2
\]

ดังนั้น

\[
\vec{R}=8\hat{i}+2\hat{j}
\]

ขนาดคือ

\[
R=\sqrt{8^2+2^2}=\sqrt{68}=2\sqrt{17}
\]

มุมกับแกน $+x$ หาได้จาก

\[
\tan\theta=\frac{2}{8}=\frac{1}{4}
\]
\end{examplebox}

\begin{examplebox}{ตัวอย่างที่ 8: ลบเวกเตอร์ด้วยองค์ประกอบ}
กำหนดให้

\[
\vec{A}=7\hat{i}+3\hat{j}
\]

และ

\[
\vec{B}=2\hat{i}+9\hat{j}
\]

หา

\[
\vec{D}=\vec{A}-\vec{B}
\]

องค์ประกอบตามแกน $x$

\[
D_x=7-2=5
\]

องค์ประกอบตามแกน $y$

\[
D_y=3-9=-6
\]

ดังนั้น

\[
\vec{D}=5\hat{i}-6\hat{j}
\]
\end{examplebox}

\begin{exercisebox}{แบบฝึกหัด 3.5}
\begin{enumerate}
    \item ให้ $\vec{A}=2\hat{i}+5\hat{j}$ และ $\vec{B}=4\hat{i}-3\hat{j}$ จงหา $\vec{A}+\vec{B}$
    \item ให้ $\vec{A}=6\hat{i}-2\hat{j}$ และ $\vec{B}=-3\hat{i}+7\hat{j}$ จงหา $\vec{A}-\vec{B}$
    \item เวกเตอร์สองตัวคือ $\vec{A}=10\hat{i}$ และ $\vec{B}=10\hat{j}$ จงหาขนาดของ $\vec{A}+\vec{B}$
    \item เวกเตอร์ลัพธ์มี $R_x=-8$ และ $R_y=6$ จงหาขนาดและควอดแรนต์ของเวกเตอร์ลัพธ์
    \item เวกเตอร์ $\vec{A}$ มีขนาด $20$ ทำมุม $30^\circ$ กับแกน $+x$ และ $\vec{B}$ มีขนาด $10$ ชี้ไปทางแกน $-y$ จงหาองค์ประกอบของ $\vec{A}+\vec{B}$
\end{enumerate}
\end{exercisebox}

% =========================================================
\section{เวกเตอร์ที่เป็นฟังก์ชันของเวลา}

ส่วนนี้เป็นการเตรียมเข้าสู่บทการเคลื่อนที่สองมิติ ถ้าตำแหน่งของอนุภาคในระนาบเขียนเป็น

\[
\vec{r}(t)=x(t)\hat{i}+y(t)\hat{j}
\]

ความเร็วคืออัตราการเปลี่ยนของเวกเตอร์ตำแหน่งตามเวลา

\[
\vec{v}(t)=\frac{d\vec{r}}{dt}
\]

เพราะ $\hat{i}$ และ $\hat{j}$ เป็นทิศคงที่ในระบบพิกัดฉาก จึงหาอนุพันธ์ทีละองค์ประกอบได้

\[
\vec{v}(t)=\frac{dx}{dt}\hat{i}+\frac{dy}{dt}\hat{j}
\]

ดังนั้น

\[
\vec{v}(t)=v_x(t)\hat{i}+v_y(t)\hat{j}
\]

โดย

\[
v_x(t)=\frac{dx}{dt}
\]

\[
v_y(t)=\frac{dy}{dt}
\]

ในทำนองเดียวกัน

\[
\vec{a}(t)=\frac{d\vec{v}}{dt}
\]

\[
\vec{a}(t)=a_x(t)\hat{i}+a_y(t)\hat{j}
\]

\begin{examplebox}{ตัวอย่างที่ 9: เวกเตอร์ตำแหน่งที่ขึ้นกับเวลา}
กำหนดให้

\[
\vec{r}(t)=(3t^2)\hat{i}+(4t)\hat{j}
\]

หาความเร็ว

\[
\vec{v}(t)=\frac{d\vec{r}}{dt}
\]

\[
\vec{v}(t)=6t\hat{i}+4\hat{j}
\]

หาความเร่ง

\[
\vec{a}(t)=\frac{d\vec{v}}{dt}
\]

\[
\vec{a}(t)=6\hat{i}
\]
\end{examplebox}

% =========================================================
\section{ชุดโจทย์รวมท้ายบท: คละระดับสำหรับ สอวน. ค่าย 1}

\subsection*{ระดับพื้นฐาน}

\begin{enumerate}
    \item จงแปลงเวกเตอร์ขนาด $10$ ทำมุม $60^\circ$ กับแกน $+x$ ให้อยู่ในรูปองค์ประกอบ
    \item จงหาขนาดของเวกเตอร์ $\vec{A}=5\hat{i}+12\hat{j}$
    \item ให้ $\vec{A}=3\hat{i}+4\hat{j}$ และ $\vec{B}=2\hat{i}-7\hat{j}$ จงหา $\vec{A}+\vec{B}$
    \item จุด $P$ มีพิกัด $(-6,8)$ จงหา $r$ และบอกควอดแรนต์
\end{enumerate}

\subsection*{ระดับกลาง}

\begin{enumerate}
    \setcounter{enumi}{4}
    \item เวกเตอร์ $\vec{A}$ มีขนาด $20$ ทำมุม $30^\circ$ กับแกน $+x$ และเวกเตอร์ $\vec{B}$ มีขนาด $10$ ทำมุม $120^\circ$ กับแกน $+x$ จงหาองค์ประกอบของ $\vec{A}+\vec{B}$
    \item แรงสองแรงขนาด $F$ เท่ากัน ทำมุม $\theta$ กับแนวดิ่งคนละด้าน จงหาแรงลัพธ์
    \item เวกเตอร์หนึ่งมีขนาด $A$ และมี $A_x=-\frac{A}{2}$ จงหามุมที่เวกเตอร์ทำกับแกน $+x$
    \item ให้ $\vec{A}=a\hat{i}+2a\hat{j}$ และ $\vec{B}=3a\hat{i}-a\hat{j}$ จงหาขนาดของ $\vec{A}-\vec{B}$
\end{enumerate}

\subsection*{ระดับยาก}

\begin{enumerate}
    \setcounter{enumi}{8}
    \item เวกเตอร์สามตัวมีผลรวมเป็นศูนย์ ถ้า
    \[
    \vec{A}=4\hat{i}+3\hat{j}
    \]
    และ
    \[
    \vec{B}=-2\hat{i}+5\hat{j}
    \]
    จงหา $\vec{C}$

    \item เวกเตอร์ $\vec{A}$ มีขนาด $A$ ทำมุม $\theta$ กับแกน $+x$ เวกเตอร์ $\vec{B}$ มีขนาด $B$ ทำมุม $-\theta$ กับแกน $+x$ จงหาองค์ประกอบของ $\vec{A}+\vec{B}$

    \item มวลหนึ่งมีโมเมนตัมเริ่มต้น $mu\hat{i}$ แล้วแตกออกเป็นสองชิ้น ชิ้นแรกมีมวล $\alpha m$ และมีความเร็ว $v\hat{j}$ จงหาทิศของความเร็วชิ้นที่สองในรูปของ $\alpha,u,v$

    \item วัตถุถูกยิงด้วยความเร็วต้น $u$ ทำมุม $\theta$ กับแนวระดับ จงเขียนเวกเตอร์ความเร็วต้นในรูปเวกเตอร์หน่วย

    \item เวกเตอร์ตำแหน่งของอนุภาคคือ
    \[
    \vec{r}(t)=(2t^2-1)\hat{i}+(3t+4)\hat{j}
    \]
    จงหา $\vec{v}(t)$ และ $\vec{a}(t)$

    \item เวกเตอร์สองตัวมีขนาดเท่ากันคือ $A$ และทำมุมระหว่างกัน $60^\circ$ จงหาขนาดของผลบวก โดยใช้การแตกองค์ประกอบ
\end{enumerate}

% =========================================================
\section{เฉลยย่อท้ายบท}

\begin{enumerate}
    \item
    \[
    A_x=10\cos60^\circ=5
    \]
    \[
    A_y=10\sin60^\circ=5\sqrt{3}
    \]

    \item
    \[
    A=\sqrt{5^2+12^2}=13
    \]

    \item
    \[
    \vec{A}+\vec{B}=(3+2)\hat{i}+(4-7)\hat{j}
    \]
    \[
    \vec{A}+\vec{B}=5\hat{i}-3\hat{j}
    \]

    \item
    \[
    r=\sqrt{(-6)^2+8^2}=10
    \]
    จุดอยู่ควอดแรนต์ที่ II

    \item
    \[
    A_x=20\cos30^\circ=10\sqrt{3}
    \]
    \[
    A_y=20\sin30^\circ=10
    \]
    \[
    B_x=10\cos120^\circ=-5
    \]
    \[
    B_y=10\sin120^\circ=5\sqrt{3}
    \]
    ดังนั้น
    \[
    R_x=10\sqrt{3}-5
    \]
    \[
    R_y=10+5\sqrt{3}
    \]

    \item
    องค์ประกอบแนวนอนหักล้างกัน และองค์ประกอบแนวดิ่งรวมกัน
    \[
    R=2F\cos\theta
    \]
    ทิศตามแนวดิ่ง

    \item
    \[
    \cos\theta=\frac{A_x}{A}=-\frac{1}{2}
    \]
    ดังนั้น
    \[
    \theta=120^\circ
    \]
    หรือถ้าอยู่ควอดแรนต์ III อาจเป็น $240^\circ$ ต้องใช้ข้อมูลของ $A_y$ เพิ่ม

    \item
    \[
    \vec{A}-\vec{B}=(a-3a)\hat{i}+(2a-(-a))\hat{j}
    \]
    \[
    \vec{A}-\vec{B}=-2a\hat{i}+3a\hat{j}
    \]
    \[
    |\vec{A}-\vec{B}|=\sqrt{(-2a)^2+(3a)^2}=\sqrt{13}a
    \]

    \item
    \[
    \vec{A}+\vec{B}+\vec{C}=0
    \]
    \[
    \vec{C}=-(\vec{A}+\vec{B})
    \]
    \[
    \vec{A}+\vec{B}=2\hat{i}+8\hat{j}
    \]
    \[
    \vec{C}=-2\hat{i}-8\hat{j}
    \]

    \item
    \[
    A_x=A\cos\theta,\quad A_y=A\sin\theta
    \]
    \[
    B_x=B\cos\theta,\quad B_y=-B\sin\theta
    \]
    ดังนั้น
    \[
    R_x=(A+B)\cos\theta
    \]
    \[
    R_y=(A-B)\sin\theta
    \]

    \item
    ใช้การอนุรักษ์โมเมนตัม
    \[
    \tan\theta=\frac{\alpha v}{u}
    \]
    โดย $\theta$ เป็นมุมที่ชิ้นที่สองทำกับแกน $+x$ ลงด้านล่าง

    \item
    \[
    \vec{u}=u\cos\theta\hat{i}+u\sin\theta\hat{j}
    \]

    \item
    \[
    \vec{v}(t)=\frac{d\vec{r}}{dt}=4t\hat{i}+3\hat{j}
    \]
    \[
    \vec{a}(t)=\frac{d\vec{v}}{dt}=4\hat{i}
    \]

    \item
    เลือกให้เวกเตอร์แรกอยู่ตามแกน $x$
    \[
    \vec{A}=A\hat{i}
    \]
    เวกเตอร์ที่สองทำมุม $60^\circ$
    \[
    \vec{B}=A\cos60^\circ\hat{i}+A\sin60^\circ\hat{j}
    \]
    ดังนั้น
    \[
    R_x=A+\frac{A}{2}=\frac{3A}{2}
    \]
    \[
    R_y=\frac{\sqrt{3}A}{2}
    \]
    \[
    R=\sqrt{\left(\frac{3A}{2}\right)^2+\left(\frac{\sqrt{3}A}{2}\right)^2}
    \]
    \[
    R=A\sqrt{3}
    \]
\end{enumerate}

% =========================================================
\section{สรุปสูตรสำคัญ}

\begin{tcolorbox}[title=สรุปบทที่ 3: เวกเตอร์]
พิกัดฉากและพิกัดเชิงขั้ว

\[
x=r\cos\theta
\]

\[
y=r\sin\theta
\]

\[
r=\sqrt{x^2+y^2}
\]

\[
\tan\theta=\frac{y}{x}
\]

องค์ประกอบของเวกเตอร์

\[
A_x=A\cos\theta
\]

\[
A_y=A\sin\theta
\]

\[
A=\sqrt{A_x^2+A_y^2}
\]

\[
\tan\theta=\frac{A_y}{A_x}
\]

เวกเตอร์หน่วย

\[
\vec{A}=A_x\hat{i}+A_y\hat{j}
\]

\[
\vec{A}=A_x\hat{i}+A_y\hat{j}+A_z\hat{k}
\]

การบวกเวกเตอร์

\[
R_x=A_x+B_x
\]

\[
R_y=A_y+B_y
\]

\[
\vec{R}=R_x\hat{i}+R_y\hat{j}
\]

\[
R=\sqrt{R_x^2+R_y^2}
\]

การลบเวกเตอร์

\[
\vec{A}-\vec{B}=\vec{A}+(-\vec{B})
\]

\[
D_x=A_x-B_x
\]

\[
D_y=A_y-B_y
\]
\end{tcolorbox}

% =========================================================
\section{ข้อสอบจริง สอวน. ที่ใช้ความรู้เรื่องเวกเตอร์}


% =========================================================
\subsection{อ่านเสริม: ปี 2566 ข้อ 3 การระเบิดและทิศของโมเมนตัม}

\vspace{1.5em}

\IfFileExists{img/posn66q3.png}{
\begin{center}
    \includegraphics[width=1\linewidth]{img/posn66q3.png}
\end{center}
}{
\begin{tcolorbox}[title=ตำแหน่งภาพที่แนะนำ]
ให้ตัดภาพข้อสอบปี 2566 ข้อ 3 แล้วบันทึกเป็น

\[
\texttt{img/posn66q3.png}
\]
\end{tcolorbox}
}

\vspace{1.5em}

\begin{tcolorbox}[title=หมายเหตุ]
ข้อนี้ใช้ความรู้เกินบทเวกเตอร์เล็กน้อย เพราะต้องใช้หลักการอนุรักษ์โมเมนตัม แต่แกนหลักของการคำนวณยังเป็นการแยกเวกเตอร์เป็นองค์ประกอบตามแกน $x$ และ $y$ จึงเหมาะใช้เป็นตัวอย่างเสริมท้ายบท
\end{tcolorbox}

\begin{examplebox}{เฉลยละเอียด}
มวล $m$ เคลื่อนที่ด้วยความเร็ว $u$ ตามแกน $+x$ แล้วระเบิดออกเป็นสองชิ้น

ชิ้นแรกมีมวล

\[
\alpha m
\]

และเคลื่อนที่ด้วยความเร็ว $v$ ตามแกน $+y$

ชิ้นที่สองมีมวล

\[
(1-\alpha)m
\]

ต้องการหาทิศของความเร็วชิ้นที่สอง

ก่อนระเบิด โมเมนตัมรวมอยู่ตามแกน $+x$ เท่านั้น

\[
\vec{p}_{\text{before}}=mu\hat{i}
\]

หลังระเบิด โมเมนตัมของชิ้นแรกคือ

\[
\vec{p}_1=\alpha mv\hat{j}
\]

ให้ความเร็วของชิ้นที่สองเป็น

\[
\vec{v}_2=v_{2x}\hat{i}+v_{2y}\hat{j}
\]

ดังนั้นโมเมนตัมของชิ้นที่สองคือ

\[
\vec{p}_2=(1-\alpha)m(v_{2x}\hat{i}+v_{2y}\hat{j})
\]

ใช้การอนุรักษ์โมเมนตัม

\[
\vec{p}_{\text{before}}=\vec{p}_1+\vec{p}_2
\]

แทนค่า

\[
mu\hat{i}
=
\alpha mv\hat{j}
+
(1-\alpha)m(v_{2x}\hat{i}+v_{2y}\hat{j})
\]

แยกแกน $x$

\[
mu=(1-\alpha)mv_{2x}
\]

\[
v_{2x}=\frac{u}{1-\alpha}
\]

แยกแกน $y$

\[
0=\alpha mv+(1-\alpha)mv_{2y}
\]

\[
v_{2y}=-\frac{\alpha v}{1-\alpha}
\]

ดังนั้นชิ้นที่สองต้องมีองค์ประกอบความเร็วลงด้านล่าง เพื่อหักล้างโมเมนตัมแกน $y$ ของชิ้นแรก

ถ้าให้ $\theta$ เป็นมุมที่ชิ้นที่สองทำกับแกน $+x$ ลงด้านล่าง

\[
\tan\theta=\frac{|v_{2y}|}{v_{2x}}
\]

แทนค่า

\[
\tan\theta
=
\frac{
\frac{\alpha v}{1-\alpha}
}{
\frac{u}{1-\alpha}
}
\]

\[
\tan\theta=\frac{\alpha v}{u}
\]

ดังนั้น

\[
\boxed{
\theta=\tan^{-1}\left(\frac{\alpha v}{u}\right)
}
\]

หรือถ้าตอบในรูปของค่าแทนเจนต์

\[
\boxed{
\tan\theta=\frac{\alpha v}{u}
}
\]
\end{examplebox}
% =========================================================
\subsection{ปี 2568 ข้อ 2: การชนยืดหยุ่นของทรงกลมผิวลื่น}

\vspace{1.5em}

\IfFileExists{img/posn68q2.png}{
\begin{center}
    \includegraphics[width=1\linewidth]{img/posn68q2.png}
\end{center}
}{
\begin{tcolorbox}[title=ตำแหน่งภาพที่แนะนำ]
ให้ตัดภาพข้อสอบปี 2568 ข้อ 2 แล้วบันทึกเป็น

\[
\texttt{img/posn68q2.png}
\]

ข้อนี้เป็นข้อเสริมท้ายบทเวกเตอร์ เพราะแกนหลักคือการแตกเวกเตอร์ตามแนวเส้นศูนย์กลางและแนวสัมผัส แต่มีแนวคิดการชนยืดหยุ่นเกินบทนี้เล็กน้อย
\end{tcolorbox}
}

\vspace{1.5em}

\begin{examplebox}{เฉลยแนวคิด}
ทรงกลม $A$ และ $B$ มีมวลเท่ากัน ผิวเกลี้ยงและลื่น ดังนั้นแรงกระทำระหว่างกันขณะชนจะอยู่ตามแนวเส้นศูนย์กลางของทรงกลมเท่านั้น

ให้แนวเส้นศูนย์กลางขณะชนทำมุม

\[
30^\circ
\]

กับทิศความเร็วเดิมของ $A$

ความเร็วต้นของ $A$ มีขนาด

\[
u
\]

แตกความเร็วของ $A$ เป็นสองแนว

\[
\text{แนวตั้งฉากกับผิวสัมผัส หรือแนวเส้นศูนย์กลาง}
\]

และ

\[
\text{แนวสัมผัส}
\]

องค์ประกอบตามแนวเส้นศูนย์กลางคือ

\[
u_n=u\cos30^\circ
\]

องค์ประกอบตามแนวสัมผัสคือ

\[
u_t=u\sin30^\circ
\]

เพราะผิวทรงกลมลื่น แรงกระทำขณะชนไม่มีองค์ประกอบตามแนวสัมผัส ดังนั้นองค์ประกอบความเร็วตามแนวสัมผัสของ $A$ ไม่เปลี่ยน

\[
v_{A,t}=u_t
\]

สำหรับการชนยืดหยุ่นของมวลเท่ากัน โดย $B$ อยู่นิ่งก่อนชน องค์ประกอบตามแนวเส้นศูนย์กลางของ $A$ จะถูกถ่ายไปให้ $B$ ส่วน $A$ จึงเหลือเฉพาะองค์ประกอบตามแนวสัมผัส

ดังนั้นขนาดของความเร็วของ $A$ หลังชนคือ

\[
v_A=u_t
\]

\[
v_A=u\sin30^\circ
\]

\[
v_A=\frac{u}{2}
\]

ดังนั้น

\[
\boxed{v_A=\frac{u}{2}}
\]
\end{examplebox}